\subsection{模的 n 次对称幂与对称多重线性映射}

设 X, Y 为两个集合，n 为整数 \(\geqslant1\) 。从 \(X^{n}\) 到 Y 的对称映射是指任意满足以下条件的映射 \(f: X^{n} \to Y\) ：对于每个置换 \(\sigma \in S_{n}\) 以及每个元素 \((x_{i}) \in X^{n}\) ，

\[
f (x _ {\sigma (1)}, x _ {\sigma (2)}, \dots , x _ {\sigma (n)}) = f (x _ {1}, x _ {2}, \dots , x _ {n}).\tag{5}
\]

由于交换两个相邻整数的对换生成群 \(\mathfrak{S}_n\) （I，§ 5，第 7 号），只需条件（5）对 \(\sigma\) 为这样一个对换的情形成立即可。

当 Y 是交换环 A 上的模时，显然从 \(X^{n}\) 到 Y 的对称映射的集合是 A-模 \(Y^{X^{n}}\) （由所有从 \(X^{n}\) 到 Y 的映射组成）的一个子模。

命题6. 设 A 是一个交换环，M 和 N 是两个 A-模。如果对每个 A-线性映射 \(g: \mathbf{S}^n(\mathbf{M}) \to \mathbf{N} (n \geqslant 1)\) 令其对应 n-线性映射

\[
\left(x _ {1}, x _ {2}, \dots , x _ {n}\right) \mapsto g \left(x _ {1} x _ {2} \dots x _ {n}\right)\tag{6}
\]

（其中右端的乘积取在代数 \(\mathbf{S}(\mathbf{M})\) 中），则得到 A-模 \(\operatorname{Hom}_{\mathbf{A}}(\mathbf{S}^{n}(\mathbf{M}), \mathbf{N})\) 到由对称 \(n\) -线性映射（从 \(\mathbf{M}^{n}\) 到 \(\mathbf{N}\) ）组成的 A-模上的一个双射 A-线性映射。

回忆（II，§3，第9号），存在从 A-模 \(\operatorname{Hom}_{\mathbf{A}}(\mathbb{T}^{n}(\mathbf{M}), \mathbf{N})\) 到 A-模 \(\mathcal{L}_{n}(\mathbf{M}, \ldots, \mathbf{M}; \mathbf{N})\) （由所有 \(n\) -线性映射（从 \(\mathbf{M}^{n}\) 到 \(\mathbf{N}\) ）组成）上的典范双射，它是通过把每个 A-线性映射 \(f: \mathbb{T}^{n}(\mathbf{M}) \to \mathbf{N}\) 对应到 \(n\) -线性映射而得到的

\[
\bar {f} \colon (x _ {1}, x _ {2}, \dots , x _ {n}) \mapsto f (x _ {1} \otimes x _ {2} \otimes \dots \otimes x _ {n}).\tag{7}
\]

另一方面，A-线性映射 \(g \colon \mathbb{S}^n(\mathbf{M}) \to \mathbf{N}\) 一一对应于满足以下条件的 A-线性映射 \(f \colon \mathbb{T}^n(\mathbf{M}) \to \mathbf{N}\) ：\(f\) 在 \(\mathfrak{J}_n'\) 上为零；对应方式是把 \(g\) 对应到映射 \(f = g \circ p_n\) ，其中

\[
p _ {n} \colon \mathbf {T} ^ {n} (\mathbf {M}) \rightarrow \mathbf {S} ^ {n} (\mathbf {M}) = \mathbf {T} ^ {n} (\mathbf {M}) / \mathfrak {J} _ {n} ^ {\prime}
\]

是典范同态（II，§ 2，第 1 号，定理 1）。但由于 \(\mathfrak{J}_n'\) 是形如

\[
\left(u _ {1} \otimes u _ {2} \otimes \dots \otimes u _ {p}\right) \otimes (x \otimes y - y \otimes x) \otimes \left(v _ {1} \otimes \dots \otimes v _ {n - p - 2}\right)
\]

\((x, y, u_{i}, v_{j} \text{ in } \mathbf{M})\) 的元素的线性组合，故说函数 f 具有形式 \(g \circ p_{n}\) 就意味着相应的 n-线性函数 \(\bar{f}\) 满足关系

\[
\bar {f} (u _ {1}, \dots , u _ {p}, x, y, v _ {1}, \dots , v _ {n - p - 2}) = \bar {f} (u _ {1}, \dots , u _ {p}, y, x, v _ {1}, \dots , v _ {n - p - 2});
\]

换言之，根据上文所见，这意味着 \(\bar{f}\) 是对称的；由此得到命题，同时考虑到以下事实：

\[
p _ {n} (x _ {1} \otimes x _ {2} \otimes \dots \otimes x _ {n}) = x _ {1} x _ {2} \dots x _ {n}
\]

（其中 \(x_{i} \in M\) ）。

A-模 \(\mathbf{S}^{n}(\mathbf{M})\) 称为 M 的第 n 次对称幂。对于每个 A-模同态 \(u: M \to N\) ，映射 \(\mathbf{S}^{n}(u): \mathbf{S}^{n}(\mathbf{M}) \to \mathbf{S}^{n}(\mathbf{N})\) （它与 \(\mathbf{S}(u)\) 在 \(\mathbf{S}^{n}(\mathbf{M})\) 上相一致）称为 u 的第 n 次对称幂。

注记。设 \(\sigma\) 是 \(\mathfrak{S}_n\) 中的一个置换；由于映射

\[
\left(x _ {1}, x _ {2}, \dots , x _ {n}\right) \mapsto x _ {\sigma^ {- 1} (1)} \otimes x _ {\sigma^ {- 1} (2)} \otimes \dots \otimes x _ {\sigma^ {- 1} (n)}
\]

（从 \(\mathbf{M}^n\) 到 \(\mathbf{T}^n (\mathbf{M})\) ）是 A-多重线性的，它可以唯一地写成

\[
(x _ {1}, \dots , x _ {n}) \mapsto u _ {\sigma} (x _ {1} \otimes x _ {2} \otimes \dots \otimes x _ {n}),
\]

（其中 \(u_{\sigma}\) 是 A-模 \(\mathsf{T}^n(\mathbf{M})\) 的一个自同态，也记作 \(z \mapsto \sigma. z\) 。显然，若 \(\sigma\) 是 \(\mathfrak{S}_n\) 的单位元，则 \(u_{\sigma}\) 是恒等映射；另一方面，写出 \(y_i = x_{\sigma^{-1}(i)}\) ，我们得到：对每个置换 \(\tau \in \mathfrak{S}_n, y_{\tau^{-1}(i)} = x_{\sigma^{-1}(\tau^{-1}(i))}\) ，因此 \(\tau. (\sigma. z) = (\tau \sigma). z\) ；换言之，A-模 \(\mathsf{T}^n(\mathbf{M})\) 是一个左 \(\mathfrak{S}_n\text{-set}\) （运算为 \((\sigma, z) \mapsto \sigma. z\) ，I，§ 5，第 1 号）。 \(\mathsf{T}^n(\mathbf{M})\) 中使得 \(\sigma. z = z\) （对所有 \(\sigma \in \mathfrak{S}_n\) ）的元素称为 \(n\) 阶（反变）对称张量；它们构成 \(\mathsf{S}_n'(\mathbf{M})\) ，即 \(\mathsf{T}^n(\mathbf{M})\) 的一个子 A-模。

对所有 \(z \in \mathbf{T}^n(\mathbf{M})\) ，我们记 \(s.z = \sum_{\sigma \in \mathfrak{S}_n} \sigma.z\) ，并称 \(s.z\) 为张量 \(z\) 的对称化；显然 \(s.z\) 是对称张量，因此 \(z \mapsto s.z\) 是 \(\mathbf{T}^n(\mathbf{M})\) 的自同态，其像 \(\mathbf{S}_n''(\mathbf{M})\) 包含于 \(\mathbf{S}_n'(\mathbf{M})\) ；一般地， \(\mathbf{S}_n''(\mathbf{M}) \neq \mathbf{S}_n'( \mathbf{M})\) （习题5）。若 \(z\) 是对称张量，则 \(s.z = n!z\) ；由此可知，当 \(n!\) 在 \(\mathbf{A}\) 中可逆时，自同态 \(z \mapsto (n!)^{-1}s.z\) 是 \(\mathbf{T}^n(\mathbf{M})\) 的投影算子（II，§1，第8号），其像为 \(\mathbf{S}_n'(\mathbf{M}) = \mathbf{S}_n''(\mathbf{M})\) ；此外，这个投影算子的核恰好是 \(\mathfrak{J}_n'\) 。因为，显然 \(\sigma(\mathfrak{J}_n') \subset \mathfrak{J}_n'\) 对所有 \(\sigma \in \mathfrak{S}_n\) 成立，且 \(\mathfrak{J}_n'\) 按定义由形如 \(z - \rho.z\) 的张量生成，其中 \(\rho\) 是交换 {[}1, n{]} 中两个相邻数字的对换；此外，若 \(\sigma, \tau\) 是 \(\mathfrak{S}_n\) 中的两个置换，则 \(z - (\sigma\tau).z = z - \sigma.z + \sigma.(z - \tau.z)\) ，由此可得（因为 \(\mathfrak{S}_n\) 中的每个置换都是交换两个相邻数字的对换之积） \(z - \sigma.z \in \mathfrak{J}_n'\) 对所有 \(z \in \mathbf{T}^n(\mathbf{M})\) 和 \(\sigma \in \mathfrak{S}_n\) 成立。因此（始终假设 \(n!\) 在 \(\mathbf{A}\) 中可逆），可见 \(z - (n!)^{-1}s.z = \sum_{\sigma \in \mathfrak{S}_n}(n!)^{-1}(z - \sigma.z) \in \mathfrak{J}_n'\) 对所有 \(z \in \mathbf{T}^n(\mathbf{M})\) 成立，这就证明了我们的断言。

当 \(n!\) 在A中可逆时，子模 \(S_n'(M)\) 和 \(\mathfrak{S}_n'\) （ \(T^n(M)\) 中的）因此是互补的，并且典范同态在 \(S_n'(M)\) 上的限制（该典范同态即 \(T^n(M) \to S^n(M) = T^n(M)/\mathfrak{J}_n'\) ）是一个 A-模同构，这使我们在所考虑的情形下可以把 \(n\) 阶对称张量与 M 的第 \(n\) 次对称幂的元素等同起来。但请注意，这种等同与乘法不相容：两个对称张量在 \(T(M)\) 中的乘积一般不是对称的，因此其在 \(S(M)\) 中的像并不是所考虑的对称张量的像的乘积。

